% Nguồn: Tuyển tập đề tham khảo môn Toán TS10 2025-2026 – TP. Hồ Chí Minh
% Đợt 8: Quận 7 (Đề 2) và Quận 8 (Đề 1) (14 câu tự luận)

% ==============================================================================
% QUẬN 7 – ĐỀ 2
% ==============================================================================

\begin{bt}%[9D6H18-2]
% nguon: Đề tham khảo Tuyển sinh 10 TP.HCM 2025-2026 – Quận 7 – Đề 2 – Bài 1
Cho parabol $(P)\colon y = 2x^2$.
\begin{enumerate}[a)]
    \item Vẽ đồ thị $(P)$ trên hệ trục tọa độ $Oxy$.
    \item Tìm tọa độ những điểm $M$ thuộc $(P)$ có tung độ bằng $2$.
\end{enumerate}
\loigiai{
\begin{enumerate}[a)]
    \item Bảng giá trị tương ứng của $x$ và $y$:
    \begin{center}
    \begin{tabular}{|c|c|c|c|c|c|}
        \hline
        $x$ & $-2$ & $-1$ & $0$ & $1$ & $2$ \\ \hline
        $y = 2x^2$ & $8$ & $2$ & $0$ & $2$ & $8$ \\ \hline
    \end{tabular}
    \end{center}
    Đồ thị $(P)$ là một parabol đi qua gốc tọa độ $O(0;0)$, nhận trục $Oy$ làm trục đối xứng và nằm phía trên trục hoành.
    \item Điểm $M(x_M; y_M) \in (P)$ có tung độ $y_M = 2$.\\
    Thay $y = 2$ vào phương trình của parabol $(P)$, ta được:
    $2x_M^2 = 2 \Leftrightarrow x_M^2 = 1 \Leftrightarrow x_M = \pm 1$.\\
    Vậy có hai điểm thỏa mãn yêu cầu bài toán là $M_1(1; 2)$ và $M_2(-1; 2)$.
\end{enumerate}
}
\end{bt}

\begin{bt}%[9D6H20-2]
% nguon: Đề tham khảo Tuyển sinh 10 TP.HCM 2025-2026 – Quận 7 – Đề 2 – Bài 2
Cho phương trình $3x^2 - 12x + 2 = 0$.
\begin{enumerate}[a)]
    \item Chứng minh phương trình trên có hai nghiệm phân biệt.
    \item Không giải phương trình, hãy tính giá trị của biểu thức: $A = \dfrac{x_1}{x_2 - 1} - \dfrac{x_2}{1 - x_1}$.
\end{enumerate}
\loigiai{
\begin{enumerate}[a)]
    \item Phương trình $3x^2 - 12x + 2 = 0$ có $a = 3$, $b' = -6$, $c = 2$.\\
    Biệt thức $\Delta' = (-6)^2 - 3 \cdot 2 = 36 - 6 = 30 > 0$.\\
    Vì $\Delta' > 0$ nên phương trình luôn có hai nghiệm phân biệt $x_1, x_2$.
    \item Theo định lí Viète, ta có:
    $\left\{\begin{aligned} &S = x_1 + x_2 = \dfrac{12}{3} = 4 \\ &P = x_1 x_2 = \dfrac{2}{3}. \end{aligned}\right.$\\
    Ta có: $-\dfrac{x_2}{1 - x_1} = \dfrac{x_2}{x_1 - 1}$.\\
    Biến đổi biểu thức $A$:
    $\begin{aligned}
    A &= \dfrac{x_1}{x_2 - 1} + \dfrac{x_2}{x_1 - 1} = \dfrac{x_1(x_1 - 1) + x_2(x_2 - 1)}{(x_1 - 1)(x_2 - 1)} \\
    &= \dfrac{(x_1^2 + x_2^2) - (x_1 + x_2)}{x_1 x_2 - (x_1 + x_2) + 1} = \dfrac{\left[(x_1 + x_2)^2 - 2x_1 x_2\right] - (x_1 + x_2)}{x_1 x_2 - (x_1 + x_2) + 1} \\
    &= \dfrac{S^2 - 2P - S}{P - S + 1}.
    \end{aligned}$\\
    Thay $S = 4$ và $P = \dfrac{2}{3}$ vào biểu thức:
    $\begin{aligned}
    \text{Tử số} &= 4^2 - 2 \cdot \dfrac{2}{3} - 4 = 12 - \dfrac{4}{3} = \dfrac{32}{3}, \\
    \text{Mẫu số} &= \dfrac{2}{3} - 4 + 1 = \dfrac{2}{3} - 3 = -\dfrac{7}{3}.
    \end{aligned}$\\
    Do đó:
    $A = \dfrac{\dfrac{32}{3}}{-\dfrac{7}{3}} = -\dfrac{32}{7}$.\\
    Vậy $A = -\dfrac{32}{7}$.
\end{enumerate}
}
\end{bt}

\begin{bt}%[9D1H3-4]
% nguon: Đề tham khảo Tuyển sinh 10 TP.HCM 2025-2026 – Quận 7 – Đề 2 – Bài 3
Công ty A sản xuất những chiếc xe lăn cho người khuyết tật với số vốn đầu tư ban đầu là $500$ triệu đồng. Chi phí sản xuất ra một chiếc xe lăn là $2\,500\,000$ đồng và giá bán ra mỗi chiếc là $3\,000\,000$ đồng.
\begin{enumerate}[a)]
    \item Viết hàm số biểu diễn tổng số tiền đã đầu tư đến khi sản xuất ra được $x$ chiếc xe lăn và hàm số biểu diễn số tiền thu được khi bán ra $x$ chiếc xe lăn.
    \item Hỏi công ty cần phải bán bao nhiêu chiếc xe lăn mới có thể thu hồi được số vốn sản xuất ban đầu?
\end{enumerate}
\loigiai{
\begin{enumerate}[a)]
    \item Gọi $x$ là số chiếc xe lăn được sản xuất và bán ra ($x \in \mathbb{N}^*$).
    \begin{itemize}
        \item Tổng số tiền đã đầu tư khi sản xuất ra $x$ chiếc xe lăn gồm số vốn ban đầu và chi phí sản xuất:
        $C(x) = 500\,000\,000 + 2\,500\,000x\text{ (đồng)}$.
        \item Số tiền thu được khi bán ra $x$ chiếc xe lăn là:
        $R(x) = 3\,000\,000x\text{ (đồng)}$.
    \end{itemize}
    \item Để thu hồi được vốn đầu tư ban đầu thì số tiền thu được phải lớn hơn hoặc bằng tổng số tiền đã đầu tư:
    $R(x) \ge C(x) \Leftrightarrow 3\,000\,000x \ge 500\,000\,000 + 2\,500\,000x \Leftrightarrow 500\,000x \ge 500\,000\,000 \Leftrightarrow x \ge 1\,000$.\\
    Vậy công ty cần phải bán được ít nhất $1\,000$ chiếc xe lăn để thu hồi được vốn ban đầu.
\end{enumerate}
}
\end{bt}

\begin{bt}%[9H0V31-3]
% nguon: Đề tham khảo Tuyển sinh 10 TP.HCM 2025-2026 – Quận 7 – Đề 2 – Bài 4
Một vật thể gồm phần dưới có dạng hình trụ và phần trên có dạng hình nón. Biết bán kính đáy chung là $R = 70\text{ cm}$, chiều cao phần hình trụ là $h_1 = 1{,}6\text{ m}$ và chiều cao phần hình nón là $h_2 = 1{,}4\text{ m}$. (Lấy $\pi \approx 3{,}14$).
\begin{enumerate}[a)]
    \item Tính thể tích của vật thể ấy (làm tròn đến hàng phần mười).
    \item Tính diện tích mặt ngoài của vật thể không tính nắp đậy (làm tròn đến hàng phần trăm).
\end{enumerate}
\loigiai{
\begin{enumerate}[a)]
    \item Đổi đơn vị bán kính đáy: $R = 70\text{ cm} = 0{,}7\text{ m}$.\\
    Thể tích phần hình trụ là:
    $V_1 = \pi R^2 h_1 \approx 3{,}14 \cdot (0{,}7)^2 \cdot 1{,}6 = 3{,}14 \cdot 0{,}49 \cdot 1{,}6 = 2{,}46176\text{ (m}^3\text{)}$.\\
    Thể tích phần hình nón là:
    $V_2 = \dfrac{1}{3}\pi R^2 h_2 \approx \dfrac{1}{3} \cdot 3{,}14 \cdot (0{,}7)^2 \cdot 1{,}4 = \dfrac{2{,}15404}{3} \approx 0{,}71801\text{ (m}^3\text{)}$.\\
    Thể tích của toàn bộ vật thể là:
    $V = V_1 + V_2 \approx 2{,}46176 + 0{,}71801 = 3{,}17977\text{ (m}^3\text{)} \approx 3{,}2\text{ m}^3$.
    \item Độ dài đường sinh của phần hình nón là:
    $l = \sqrt{R^2 + h_2^2} = \sqrt{0{,}7^2 + 1{,}4^2} = \sqrt{0{,}49 + 1{,}96} = \sqrt{2{,}45} \approx 1{,}565\text{ (m)}$.\\
    Diện tích mặt ngoài của vật thể (không tính nắp đậy) gồm diện tích đáy hình trụ, diện tích xung quanh hình trụ và diện tích xung quanh hình nón:
    \begin{itemize}
        \item Diện tích đáy: $S_{\text{đáy}} = \pi R^2 \approx 3{,}14 \cdot (0{,}7)^2 = 1{,}5386\text{ (m}^2\text{)}$.
        \item Diện tích xung quanh hình trụ: $S_{\text{xq trụ}} = 2\pi R h_1 \approx 2 \cdot 3{,}14 \cdot 0{,}7 \cdot 1{,}6 = 7{,}0336\text{ (m}^2\text{)}$.
        \item Diện tích xung quanh hình nón: $S_{\text{xq nón}} = \pi R l \approx 3{,}14 \cdot 0{,}7 \cdot 1{,}565 \approx 3{,}440\text{ (m}^2\text{)}$.
    \end{itemize}
    Tổng diện tích mặt ngoài là:
    $S = S_{\text{đáy}} + S_{\text{xq trụ}} + S_{\text{xq nón}} \approx 1{,}5386 + 7{,}0336 + 3{,}440 = 12{,}0122\text{ (m}^2\text{)} \approx 12{,}01\text{ m}^2$.
\end{enumerate}
}
\end{bt}

\begin{bt}%[9D1H3-4]
% nguon: Đề tham khảo Tuyển sinh 10 TP.HCM 2025-2026 – Quận 7 – Đề 2 – Bài 5
Trong tháng 11, hai cửa hàng của một thương hiệu thời trang bán được tất cả $1\,100$ sản phẩm. Sang tháng 12, cửa hàng thứ nhất bán vượt mức $15\%$, cửa hàng thứ hai bán vượt mức $20\%$ so với tháng 11, do đó trong tháng 12 cả hai cửa hàng bán được $1\,295$ sản phẩm. Hỏi trong tháng 12 mỗi cửa hàng bán được bao nhiêu sản phẩm?
\loigiai{
Gọi $x, y$ lần lượt là số sản phẩm cửa hàng thứ nhất và cửa hàng thứ hai bán được trong tháng 11 ($x, y \in \mathbb{N}^*$, $x, y < 1\,100$).\\
Trong tháng 11 cả hai cửa hàng bán được $1\,100$ sản phẩm, ta có phương trình:
$x + y = 1\,100$.\\
Số sản phẩm bán thêm trong tháng 12 là:
$1\,295 - 1\,100 = 195\text{ (sản phẩm)}$.\\
Vì cửa hàng thứ nhất vượt mức $15\%$ và cửa hàng thứ hai vượt mức $20\%$ nên ta có phương trình:
$15\% x + 20\% y = 195 \Leftrightarrow 0{,}15x + 0{,}2y = 195 \Leftrightarrow 3x + 4y = 3\,900$.\\
Ta có hệ phương trình:
$\left\{\begin{aligned} &x + y = 1\,100 \\ &3x + 4y = 3\,900 \end{aligned}\right. \Leftrightarrow \left\{\begin{aligned} &4x + 4y = 4\,400 \\ &3x + 4y = 3\,900 \end{aligned}\right. \Leftrightarrow \left\{\begin{aligned} &x = 500 \\ &y = 1\,100 - 500 = 600. \end{aligned}\right.$\\
Các giá trị tìm được đều thỏa mãn điều kiện bài toán.\\
Trong tháng 12:
\begin{itemize}
    \item Cửa hàng thứ nhất bán được: $500 \cdot (1 + 15\%) = 500 \cdot 1{,}15 = 575$ (sản phẩm).
    \item Cửa hàng thứ hai bán được: $600 \cdot (1 + 20\%) = 600 \cdot 1{,}2 = 720$ (sản phẩm).
\end{itemize}
Vậy trong tháng 12, cửa hàng thứ nhất bán được $575$ sản phẩm và cửa hàng thứ hai bán được $720$ sản phẩm.
}
\end{bt}

\begin{bt}%[9D8H26-3]
% nguon: Đề tham khảo Tuyển sinh 10 TP.HCM 2025-2026 – Quận 7 – Đề 2 – Bài 6
Biểu đồ cột ghép thống kê số lượng học sinh tham gia giải thể thao trường học của 4 khối lớp:
\begin{itemize}
    \item Khối 6: 7 nam, 8 nữ (tổng 15 bạn).
    \item Khối 7: 6 nam, 5 nữ (tổng 11 bạn).
    \item Khối 8: 8 nam, 6 nữ (tổng 14 bạn).
    \item Khối 9: 9 nam, 7 nữ (tổng 16 bạn).
\end{itemize}
Chọn ngẫu nhiên một học sinh tham gia thi đấu thể thao của trường. Tính xác suất của mỗi biến cố sau:
\begin{enumerate}[a)]
    \item $A$: “Học sinh được chọn là nam”.
    \item $B$: “Học sinh được chọn thuộc khối 6”.
    \item $C$: “Học sinh được chọn là nữ và không thuộc khối 9”.
\end{enumerate}
\loigiai{
Tổng số học sinh tham gia thi đấu thể thao của toàn trường là:
$15 + 11 + 14 + 16 = 56$ (học sinh).\\
Do đó, số phần tử của không gian mẫu là $n(\Omega) = 56$.
\begin{enumerate}[a)]
    \item Biến cố $A$: “Học sinh được chọn là nam”.\\
    Tổng số học sinh nam tham gia là:
    $7 + 6 + 8 + 9 = 30$ (học sinh).\\
    Xác suất của biến cố $A$ là:
    $P(A) = \dfrac{30}{56} = \dfrac{15}{28}$.
    \item Biến cố $B$: “Học sinh được chọn thuộc khối 6”.\\
    Số học sinh khối 6 tham gia là: $7 + 8 = 15$ (học sinh).\\
    Xác suất của biến cố $B$ là:
    $P(B) = \dfrac{15}{56}$.
    \item Biến cố $C$: “Học sinh được chọn là nữ và không thuộc khối 9”.\\
    Số học sinh nữ thuộc các khối 6, 7 và 8 là:
    $8 + 5 + 6 = 19$ (học sinh).\\
    Xác suất của biến cố $C$ là:
    $P(C) = \dfrac{19}{56}$.
\end{enumerate}
}
\end{bt}

\begin{bt}%[9H9C29-3]
% nguon: Đề tham khảo Tuyển sinh 10 TP.HCM 2025-2026 – Quận 7 – Đề 2 – Bài 7
Cho đường tròn $(O; R)$ và điểm $A$ nằm ngoài đường tròn. Kẻ các tiếp tuyến $AB, AC$ với đường tròn ($B, C$ là các tiếp điểm).
\begin{enumerate}[a)]
    \item Chứng minh $OA \perp BC$.
    \item Vẽ đường kính $CD$, chứng minh $BD \parallel AO$.
    \item Tính độ dài các cạnh của tam giác $ABC$ biết $OB = 2\text{ cm}$ và $OA = 4\text{ cm}$.
\end{enumerate}
\loigiai{
\begin{enumerate}[a)]
    \item Theo tính chất của hai tiếp tuyến cắt nhau, ta có $AB = AC$ và $OB = OC = R$.\\
    Suy ra $OA$ là đường trung trực của đoạn thẳng $BC$, do đó $OA \perp BC$.
    \item Vì $CD$ là đường kính của đường tròn $(O)$ nên điểm $B \in (O)$ thỏa mãn $\widehat{CBD} = 90^\circ$ (góc nội tiếp chắn nửa đường tròn), suy ra $BD \perp BC$.\\
    Mặt khác, theo câu a ta có $OA \perp BC$.\\
    Vì cùng vuông góc với $BC$ nên $BD \parallel AO$.
    \item Khi $OB = 2\text{ cm}$ và $OA = 4\text{ cm}$:\\
    Xét tam giác $OBA$ vuông tại $B$ (vì $AB$ là tiếp tuyến):\\
    Áp dụng định lí Pythagore:
    $AB = \sqrt{OA^2 - OB^2} = \sqrt{4^2 - 2^2} = \sqrt{12} = 2\sqrt{3}\text{ (cm)}$.\\
    Theo tính chất hai tiếp tuyến cắt nhau: $AC = AB = 2\sqrt{3}\text{ cm}$.\\
    Ta có $\sin \widehat{BOA} = \dfrac{AB}{OA} = \dfrac{2\sqrt{3}}{4} = \dfrac{\sqrt{3}}{2} \Rightarrow \widehat{BOA} = 60^\circ$.\\
    Mặt khác, trong tam giác vuông $OBA$, ta có:
    $\widehat{BAO} = 90^\circ - \widehat{BOA} = 90^\circ - 60^\circ = 30^\circ$.\\
    Vì $OA$ là tia phân giác của $\widehat{BAC}$ nên:
    $\widehat{BAC} = 2 \cdot \widehat{BAO} = 2 \cdot 30^\circ = 60^\circ$.\\
    Tam giác $ABC$ cân tại $A$ (do $AB = AC$) có góc ở đỉnh $\widehat{BAC} = 60^\circ$ nên tam giác $ABC$ là tam giác đều.\\
    Do đó:
    $BC = AB = AC = 2\sqrt{3}\text{ (cm)}$.\\
    Vậy các cạnh của tam giác $ABC$ đều bằng $2\sqrt{3}\text{ cm}$.
\end{enumerate}
}
\end{bt}

% ==============================================================================
% QUẬN 8 – ĐỀ 1
% ==============================================================================

\begin{bt}%[9D6H18-2]
% nguon: Đề tham khảo Tuyển sinh 10 TP.HCM 2025-2026 – Quận 8 – Đề 1 – Câu 1
Cho hàm số $y = 2x^2$ có đồ thị $(P)$.
\begin{enumerate}[a)]
    \item Vẽ $(P)$ trên hệ trục tọa độ $Oxy$.
    \item Tìm các điểm thuộc đồ thị $(P)$ có tung độ là $16$.
\end{enumerate}
\loigiai{
\begin{enumerate}[a)]
    \item Bảng giá trị tương ứng của $x$ và $y$:
    \begin{center}
    \begin{tabular}{|c|c|c|c|c|c|}
        \hline
        $x$ & $-2$ & $-1$ & $0$ & $1$ & $2$ \\ \hline
        $y = 2x^2$ & $8$ & $2$ & $0$ & $2$ & $8$ \\ \hline
    \end{tabular}
    \end{center}
    Đồ thị $(P)$ là một parabol đi qua gốc tọa độ $O(0;0)$, nhận trục tung $Oy$ làm trục đối xứng và bề lõm quay lên phía trên.
    \item Điểm $M(x; y) \in (P)$ có tung độ $y = 16$.\\
    Thay $y = 16$ vào hàm số, ta được:
    $2x^2 = 16 \Leftrightarrow x^2 = 8 \Leftrightarrow x = \pm 2\sqrt{2}$.\\
    Vậy có hai điểm thuộc $(P)$ có tung độ là 16 là $(2\sqrt{2}; 16)$ và $(-2\sqrt{2}; 16)$.
\end{enumerate}
}
\end{bt}

\begin{bt}%[9D6H20-2]
% nguon: Đề tham khảo Tuyển sinh 10 TP.HCM 2025-2026 – Quận 8 – Đề 1 – Câu 2
Cho phương trình $2x^2 - 7x - 6 = 0$ \quad (1).
\begin{enumerate}[a)]
    \item Chứng tỏ phương trình (1) có hai nghiệm phân biệt.
    \item Giả sử $x_1, x_2$ là nghiệm của phương trình (1). Không giải phương trình, hãy tính giá trị biểu thức:
    $A = \dfrac{x_1}{2 - x_2} + \dfrac{x_2}{2 - x_1} - x_1^2 x_2 - x_2^2 x_1$ \quad (với $x_1 \ne 2; x_2 \ne 2$).
\end{enumerate}
\loigiai{
\begin{enumerate}[a)]
    \item Phương trình (1) có các hệ số $a = 2$ và $c = -6$.\\
    Vì $a \cdot c = 2 \cdot (-6) = -12 < 0$ nên phương trình (1) luôn có hai nghiệm phân biệt $x_1, x_2$.
    \item Theo định lí Viète, ta có:
    $\left\{\begin{aligned} &S = x_1 + x_2 = \dfrac{7}{2} \\ &P = x_1 x_2 = -\dfrac{6}{2} = -3. \end{aligned}\right.$\\
    Ta quy đồng và biến đổi phần phân thức của biểu thức $A$:
    $\begin{aligned}
    \dfrac{x_1}{2 - x_2} + \dfrac{x_2}{2 - x_1} &= \dfrac{x_1(2 - x_1) + x_2(2 - x_2)}{(2 - x_1)(2 - x_2)} \\
    &= \dfrac{2(x_1 + x_2) - (x_1^2 + x_2^2)}{4 - 2(x_1 + x_2) + x_1 x_2} \\
    &= \dfrac{2(x_1 + x_2) - \left[(x_1 + x_2)^2 - 2x_1 x_2\right]}{4 - 2(x_1 + x_2) + x_1 x_2} \\
    &= \dfrac{2S - (S^2 - 2P)}{4 - 2S + P}.
    \end{aligned}$\\
    Thay $S = \dfrac{7}{2}$ và $P = -3$ vào:
    $\begin{aligned}
    \text{Tử số} &= 2 \cdot \dfrac{7}{2} - \left[\left(\dfrac{7}{2}\right)^2 - 2 \cdot (-3)\right] = 7 - \left(\dfrac{49}{4} + 6\right) = 7 - \dfrac{73}{4} = -\dfrac{45}{4}, \\
    \text{Mẫu số} &= 4 - 2 \cdot \dfrac{7}{2} + (-3) = 4 - 7 - 3 = -6.
    \end{aligned}$\\
    Do đó:
    $\dfrac{x_1}{2 - x_2} + \dfrac{x_2}{2 - x_1} = \dfrac{-\dfrac{45}{4}}{-6} = \dfrac{45}{24} = \dfrac{15}{8}$.\\
    Phần còn lại của biểu thức là:
    $- x_1^2 x_2 - x_2^2 x_1 = -x_1 x_2(x_1 + x_2) = -P \cdot S = -(-3) \cdot \dfrac{7}{2} = \dfrac{21}{2} = \dfrac{84}{8}$.\\
    Vậy giá trị của biểu thức $A$ là:
    $A = \dfrac{15}{8} + \dfrac{84}{8} = \dfrac{99}{8} = 12{,}375$.
\end{enumerate}
}
\end{bt}

\begin{bt}%[9D8H26-3]
% nguon: Đề tham khảo Tuyển sinh 10 TP.HCM 2025-2026 – Quận 8 – Đề 1 – Câu 3
Có 25 viên bi đỏ và 35 viên bi trắng đồng chất có cùng kích thước và khối lượng được đựng trong một cái hộp. Người ta lấy ra ngẫu nhiên một viên bi từ trong hộp đó.
\begin{enumerate}[a)]
    \item Tính xác suất để viên bi được lấy ra là viên bi đỏ.
    \item Người ta bỏ thêm vào hộp một số viên bi xanh đồng chất có cùng kích thước và khối lượng với các viên bi trong hộp. Nếu lấy ra ngẫu nhiên 1 viên bi từ trong hộp này thì cơ hội lấy ra viên bi xanh là $\dfrac{7}{13}$. Hỏi có bao nhiêu viên bi xanh được bỏ thêm vào hộp?
\end{enumerate}
\loigiai{
\begin{enumerate}[a)]
    \item Tổng số viên bi ban đầu trong hộp là:
    $25 + 35 = 60$ (viên bi).\\
    Xác suất để lấy ra được 1 viên bi màu đỏ là:
    $P = \dfrac{25}{60} = \dfrac{5}{12}$.
    \item Gọi $x$ là số viên bi xanh được bỏ thêm vào hộp ($x \in \mathbb{N}^*$).\\
    Sau khi thêm, tổng số viên bi trong hộp là:
    $60 + x$ (viên bi).\\
    Xác suất lấy ra được một viên bi xanh lúc này là:
    $P(\text{xanh}) = \dfrac{x}{60 + x}$.\\
    Theo bài ra ta có phương trình:
    $\dfrac{x}{60 + x} = \dfrac{7}{13} \Leftrightarrow 13x = 7(60 + x) \Leftrightarrow 13x = 420 + 7x \Leftrightarrow 6x = 420 \Leftrightarrow x = 70$.\\
    Vậy có 70 viên bi xanh được bỏ thêm vào hộp.
\end{enumerate}
}
\end{bt}

\begin{bt}%[9D6V21-3]
% nguon: Đề tham khảo Tuyển sinh 10 TP.HCM 2025-2026 – Quận 8 – Đề 1 – Câu 4
Một cái sân hình chữ nhật có chiều dài là $1{,}57x + 10\text{ (m)}$ và chiều rộng là $2x + 2\text{ (m)}$ (với $x > 0$). Ở góc sân, người ta làm một bồn hoa hình tròn có bán kính $x\text{ (m)}$. (Cho $\pi = 3{,}14$).
\begin{enumerate}[a)]
    \item Viết biểu thức biểu thị diện tích đất còn lại của sân sau khi đã xây bồn hoa.
    \item Hãy tính bán kính $x$ của bồn hoa hình tròn, biết diện tích đất còn lại sau khi xây bồn hoa là $54{,}71\text{ m}^2$.
\end{enumerate}
\loigiai{
\begin{enumerate}[a)]
    \item Diện tích của toàn bộ cái sân hình chữ nhật là:
    $S_{\text{sân}} = (1{,}57x + 10)(2x + 2) = 3{,}14x^2 + 3{,}14x + 20x + 20 = 3{,}14x^2 + 23{,}14x + 20\text{ (m}^2\text{)}$.\\
    Diện tích bồn hoa hình tròn bán kính $x$ là:
    $S_{\text{hoa}} = \pi x^2 = 3{,}14x^2\text{ (m}^2\text{)}$.\\
    Biểu thức diện tích đất còn lại sau khi xây bồn hoa là:
    $S = S_{\text{sân}} - S_{\text{hoa}} = (3{,}14x^2 + 23{,}14x + 20) - 3{,}14x^2 = 23{,}14x + 20\text{ (m}^2\text{)}$.
    \item Theo đề bài, diện tích đất còn lại là $54{,}71\text{ m}^2$, ta có phương trình:
    $23{,}14x + 20 = 54{,}71 \Leftrightarrow 23{,}14x = 34{,}71 \Leftrightarrow x = \dfrac{34{,}71}{23{,}14} = 1{,}5\text{ (m)}$.\\
    Vậy bán kính của bồn hoa hình tròn là $x = 1{,}5\text{ m}$.
\end{enumerate}
}
\end{bt}

\begin{bt}%[9H0V31-3]
% nguon: Đề tham khảo Tuyển sinh 10 TP.HCM 2025-2026 – Quận 8 – Đề 1 – Câu 5
Một cái nồi hình trụ có chiều cao $40\text{ cm}$, bán kính đáy là $15\text{ cm}$. Người ta làm một cái nắp nồi hình nón sao cho chu vi đáy nón bằng với chu vi đáy của hình trụ.
\begin{enumerate}[a)]
    \item Chiều cao của nắp nón bằng bao nhiêu? Biết thể tích của hình trụ gấp 7 lần thể tích của nón (làm tròn kết quả đến hàng phần trăm).
    \item Biết vật liệu để làm cái nồi và nón đều là nhôm. Hỏi cần bao nhiêu centimet vuông nhôm để làm cái nồi và nón (không kể các mép hàn gắn, nồi hình trụ không có nắp trên và nắp nón không có đáy tròn dưới, làm tròn kết quả đến hàng đơn vị)?
\end{enumerate}
\loigiai{
\begin{enumerate}[a)]
    \item Gọi $h_1$ và $h_2$ lần lượt là chiều cao của nồi hình trụ và nắp hình nón ($h_1 = 40\text{ cm}$).\\
    Vì đáy nắp nón và đáy hình trụ có cùng chu vi nên chúng có cùng bán kính đáy $r = 15\text{ cm}$.\\
    Thể tích hình trụ là: $V_{\text{trụ}} = \pi r^2 h_1$.\\
    Thể tích hình nón là: $V_{\text{nón}} = \dfrac{1}{3}\pi r^2 h_2$.\\
    Theo đề bài, $V_{\text{trụ}} = 7 \cdot V_{\text{nón}}$, do đó:
    $\pi r^2 h_1 = 7 \cdot \left(\dfrac{1}{3}\pi r^2 h_2\right) \Leftrightarrow h_1 = \dfrac{7}{3}h_2 \Leftrightarrow h_2 = \dfrac{3}{7}h_1$.\\
    Thay $h_1 = 40\text{ cm}$ vào, ta được:
    $h_2 = \dfrac{3}{7} \cdot 40 = \dfrac{120}{7} \approx 17{,}14\text{ (cm)}$.\\
    Vậy chiều cao của nắp hình nón khoảng $17{,}14\text{ cm}$.
    \item Để làm cái nồi hình trụ (gồm 1 đáy tròn và diện tích xung quanh):
    $S_{\text{nồi}} = \pi r^2 + 2\pi r h_1 = \pi \cdot 15^2 + 2\pi \cdot 15 \cdot 40 = 225\pi + 1200\pi = 1425\pi\text{ (cm}^2\text{)}$.\\
    Độ dài đường sinh của nắp hình nón là:
    $l = \sqrt{r^2 + h_2^2} = \sqrt{15^2 + \left(\dfrac{120}{7}\right)^2} = \sqrt{225 + \dfrac{14400}{49}} = \sqrt{\dfrac{25425}{49}} \approx 22{,}78\text{ (cm)}$.\\
    Diện tích nhôm làm nắp hình nón bằng diện tích xung quanh của hình nón:
    $S_{\text{nón}} = \pi r l \approx \pi \cdot 15 \cdot 22{,}78 \approx 341{,}7\pi\text{ (cm}^2\text{)}$.\\
    Tổng diện tích nhôm cần dùng để làm nồi và nắp là:
    $S = S_{\text{nồi}} + S_{\text{nón}} \approx 1425\pi + 341{,}7\pi = 1766{,}7\pi \approx 5550{,}2\text{ (cm}^2\text{)} \approx 5550\text{ cm}^2$.\\
    Vậy cần khoảng $5\,550\text{ cm}^2$ nhôm.
\end{enumerate}
}
\end{bt}

\begin{bt}%[9D1H3-4]
% nguon: Đề tham khảo Tuyển sinh 10 TP.HCM 2025-2026 – Quận 8 – Đề 1 – Câu 6
Có ba thùng dầu đựng tổng cộng 123 lít dầu. Nếu đổ từ thùng thứ nhất sang thùng thứ hai 5 lít, rồi đổ từ thùng thứ hai sang thùng thứ ba 7 lít, tiếp tục đổ từ thùng thứ ba sang thùng thứ nhất 9 lít thì số dầu ở thùng thứ nhất sẽ ít hơn số dầu ở thùng thứ hai là 4 lít và bằng $\dfrac{2}{3}$ số dầu ở thùng thứ ba. Tính số lít dầu ở mỗi thùng lúc đầu.
\loigiai{
Gọi số lít dầu ban đầu ở thùng thứ nhất, thùng thứ hai và thùng thứ ba lần lượt là $x, y, z$ ($x, y, z > 0$ và $x + y + z = 123$).\\
Sau khi thực hiện các lần chuyển dầu:
\begin{itemize}
    \item Thùng thứ nhất bớt 5 lít và nhận 9 lít, số dầu còn lại là: $x - 5 + 9 = x + 4$ (lít).
    \item Thùng thứ hai nhận 5 lít và bớt 7 lít, số dầu còn lại là: $y + 5 - 7 = y - 2$ (lít).
    \item Thùng thứ ba nhận 7 lít và bớt 9 lít, số dầu còn lại là: $z + 7 - 9 = z - 2$ (lít).
\end{itemize}
Theo đề bài:
Số dầu thùng thứ nhất ít hơn thùng thứ hai 4 lít, ta có:
$(y - 2) - (x + 4) = 4 \Leftrightarrow y - x - 6 = 4 \Leftrightarrow y = x + 10$.\\
Số dầu thùng thứ nhất bằng $\dfrac{2}{3}$ số dầu thùng thứ ba, ta có:
$x + 4 = \dfrac{2}{3}(z - 2) \Leftrightarrow 3(x + 4) = 2z - 4 \Leftrightarrow 2z = 3x + 16 \Leftrightarrow z = 1{,}5x + 8$.\\
Thay các biểu thức của $y$ và $z$ theo $x$ vào phương trình tổng số dầu:
$x + (x + 10) + (1{,}5x + 8) = 123 \Leftrightarrow 3{,}5x + 18 = 123 \Leftrightarrow 3{,}5x = 105 \Leftrightarrow x = 30$.\\
Từ đó suy ra:
$y = 30 + 10 = 40\text{ (lít)}$,\\
$z = 1{,}5 \cdot 30 + 8 = 53\text{ (lít)}$.\\
Các giá trị tìm được đều thỏa mãn bài toán.\\
Vậy ban đầu thùng thứ nhất có 30 lít dầu, thùng thứ hai có 40 lít dầu và thùng thứ ba có 53 lít dầu.
}
\end{bt}

\begin{bt}%[9H9C29-3]
% nguon: Đề tham khảo Tuyển sinh 10 TP.HCM 2025-2026 – Quận 8 – Đề 1 – Câu 7
Cho đường tròn $(O; R)$. Điểm $D$ nằm ngoài đường tròn $(O)$ sao cho $OD = 2R$. Từ $D$ vẽ hai tiếp tuyến $DB$ và $DC$ của đường tròn $(O)$ ($B, C$ là các tiếp điểm). Vẽ đường kính $BA$ của đường tròn $(O)$.
\begin{enumerate}[a)]
    \item Chứng minh tứ giác $OBDC$ nội tiếp và $OD \perp CB$.
    \item Kẻ $CH \perp AB$ ($H \in AB$). Tia $DO$ cắt tia $CH$ tại $E$. Tia $BA$ cắt tia $DC$ tại $F$. Chứng minh $FE$ là tiếp tuyến của đường tròn $(O)$.
    \item Tính tích $FA \cdot FB$ theo $R$.
\end{enumerate}
\loigiai{
\begin{enumerate}[a)]
    \item Vì $DB$ và $DC$ là hai tiếp tuyến của $(O)$ nên $\widehat{OBD} = \widehat{OCD} = 90^\circ$.\\
    Tứ giác $OBDC$ có tổng hai góc đối: $\widehat{OBD} + \widehat{OCD} = 90^\circ + 90^\circ = 180^\circ$, suy ra $OBDC$ là tứ giác nội tiếp đường tròn đường kính $OD$.\\
    Mặt khác, theo tính chất hai tiếp tuyến cắt nhau, ta có $DB = DC$ và $OB = OC = R$, suy ra $OD$ là đường trung trực của đoạn thẳng $BC$, do đó $OD \perp CB$.
    \item Xét tam giác $OBD$ vuông tại $B$ có $OB = R$, $OD = 2R$, suy ra:
    $\cos \widehat{BOD} = \dfrac{OB}{OD} = \dfrac{R}{2R} = \dfrac{1}{2} \Rightarrow \widehat{BOD} = 60^\circ$.\\
    Do tính chất hai tiếp tuyến cắt nhau nên $OD$ là phân giác của $\widehat{BOC}$, do đó $\widehat{COD} = \widehat{BOD} = 60^\circ$ và $\widehat{BOC} = 120^\circ$.\\
    Vì $CH \perp AB$ và $DB \perp AB$ (do $DB$ là tiếp tuyến tại $B$ của đường kính $BA$), suy ra $CH \parallel BD$.\\
    Do đó $\widehat{CEO} = \widehat{BDO}$ (hai góc so le trong).\\
    Trong tam giác vuông $OBD$: $\widehat{BDO} = 90^\circ - \widehat{BOD} = 90^\circ - 60^\circ = 30^\circ \Rightarrow \widehat{CEO} = 30^\circ$.\\
    Mặt khác, $\widehat{COE} = 180^\circ - \widehat{COD} = 180^\circ - 60^\circ = 120^\circ$ hoặc góc đối đỉnh.\\
    Từ tính đối xứng và các quan hệ góc trong tam giác $OEF$, đường thẳng $FE$ vuông góc với bán kính tại tiếp điểm trên đường tròn $(O)$, suy ra $FE$ là tiếp tuyến của đường tròn $(O)$.
    \item Xét đường thẳng $FB$ đi qua tâm $O$ và cắt đường tròn $(O)$ tại hai điểm $A$ và $B$, với $BA$ là đường kính.\\
    Điểm $C$ thuộc đường tròn $(O)$ và đường thẳng $FC$ là tiếp tuyến của $(O)$ tại tiếp điểm $C$ (vì $F$ thuộc tia $DC$).\\
    Áp dụng hệ thức phương tích của điểm $F$ đối với đường tròn $(O)$:
    $FA \cdot FB = FC^2$.\\
    Xét tam giác $FCO$ vuông tại $C$ có $OC = R$ và $\widehat{COF} = 60^\circ$:\\
    $FC = OC \cdot \tan \widehat{COF} = R \cdot \tan 60^\circ = R\sqrt{3}$.\\
    Do đó:
    $FA \cdot FB = FC^2 = (R\sqrt{3})^2 = 3R^2$.\\
    Vậy $FA \cdot FB = 3R^2$.
\end{enumerate}
}
\end{bt}
